\documentclass[12pt,a4paper,oneside]{article}
\usepackage[brazil]{babel}
\usepackage[utf8]{inputenc}
\usepackage[dvipsnames]{xcolor}
\usepackage{listings}
\usepackage{wrapfig}
\usepackage{olymp}

\setcounter{problem}{1}

\begin{document}

\begin{problem}{Strings numéricas}{stringnum}{2}

\begingroup
\setlength{\intextsep}{0pt}%
%\setlength{\columnsep}{0pt}

Para calcular as despesas na realização de um evento, um grupo de estudantes preparou uma planilha para anotar os valores recebidos e gastos. Infelizmente alguém anotou outros dados na planilha e agora precisam separar os valores.

Faça um programa para escrever ler os dados e escrever apenas os que representam números inteiros (isto é, que contêm apenas dígitos).

% \Notes

\InputFile

A entrada começa com um linha contendo um número inteiro $N$, seguida por $N$ strings, uma em cada linha.

Restrições: $1 \leq N \leq 100$; as strings contém apenas letras e dígitos, sem espaço em branco, e nenhuma possui mais de $10$ caracteres.

\endgroup

\OutputFile

Escreva, na ordem em que aparecem, todas as strings da entrada que são números inteiros.

% \Notes

\Example

\begin{example}
\exmp{
2
1234
PAGO
}{
1234
}%
\end{example}

\begin{example}
\exmp{
5
100
Dia10
Credito
8
7mais5
}{
100
8
}%
\end{example}

\begin{example}
\exmp{
3
um
dois
3
}{
3
}%
\end{example}

%Comentários sobre o exemplo, se necessário.

\end{problem}

\end{document}